% ch01: 原文 GLDraft_D.tex 行 212–290；§1 Introduction
% 只译散文（标题、正文、图题、脚注）；公式、\label、\ref、\eqref、\cite、图文件名保持不变。

\section{引言}\label{intro}


由 Beilinson 与 Drinfeld 最初表述的几何朗兰兹对偶（geometric Langlands duality）\cite{BD}，是把黎曼曲面 $C$ 上场的模空间所关联的范畴联系起来的一种关系。许多参与表述和分析这一对偶的要素，在量子场论中都是熟悉的。特别地，二维共形场论在其中扮演突出角色，参见 \cite{F} 的综述。若用 $\N=4$ 超杨–米尔斯理论（super Yang--Mills theory）的扭化版本（twisted version）来表述，几何朗兰兹对偶与量子场论的联系就能被更充分地理解 \cite{KW}。从这个观点看，几何朗兰兹对偶可由电–磁对偶（electric--magnetic duality）推出：一边是紧致规范群 $G$ 的 $\N=4$ 理论，另一边是基于朗兰兹对偶群或 GNO 对偶群 $G^\vee$ 的同一理论（这两个群的复化分别记作 $G_\C$ 与 $G^\vee_\C$）。\zy{朗兰兹对偶群（Langlands dual group），在物理文献中也称 GNO 对偶群，得名于 Goddard--Nuyts--Olive。对紧致连通单连通单群，$G^\vee$ 的根系是 $G$ 的对偶根系；一般情形还需处理中心与连通分支。参见 R. Langlands, ``Problems in the theory of automorphic forms'' (1970)；P. Goddard, J. Nuyts and D. Olive, ``Gauge theories and magnetic charge,'' \emph{Nucl. Phys.} B125 (1977) 1。} 对 $\N=4$ 超杨–米尔斯理论作扭化，会得到一个四维拓扑场论；按通常的构造 \cite{BaDo,L1,Fr}，它给四维流形配一个数（配分函数，partition function），给三维流形配一个向量空间（物理态空间），给二维流形配一个膜（brane）或边界条件的范畴。\footnote{\label{partial}
这一理想化描述没有计入如下事实：复李群 $G_\C$、$G^\vee_\C$ 以及相应的模空间都不是紧的。由于这种非紧性，人们很可能只能得到部分的拓扑场论。膜与物理态空间可以定义，但那些形式上会用来定义四维流形配分函数的积分，是否能真正令人满意地定义，并不清楚。（四维流形的 Donaldson 理论中会出现某种类似情形，尽管细节相当不同：那里也得不到完整的拓扑场论，因为配分函数不能对所有四维流形适当地定义。）非紧性还意味着，要定义膜与物理态空间，必须指定膜或波函数所允许的渐近行为。每一种（合理的）选择都会导致对偶的不同版本。可能的选择包括 Betti 版本与 de Rham 版本 \cite{BN}。}
通常的几何朗兰兹对偶，是两个二维流形所关联的范畴之间的对偶。

\begin{yizhu}[译者按：这里的“扭化”指什么]
这里所说的“扭化”，是指 Kapustin 与 Witten 对 $\N=4$ 超杨–米尔斯理论所作的拓扑扭化（常称 GL 扭化）：取超对称生成元的一个适当组合作为标量超荷 $Q$，并将洛伦兹群与 $R$ 对称性混合，使该超荷成为标量；扭化后应力张量成为 $Q$-恰当量，从而得到拓扑场论。原文紧接着说的“部分拓扑场论”，是指由于模空间非紧，四维配分函数未必能严格定义；但膜与物理态空间仍可定义。参见 \cite{KW} 第 2 节以及 \cite{GuW} 的讨论。
\end{yizhu}

最近，Etingof、Frenkel 与 Kazhdan \cite{EFK,EFK2,EFK3} 发现了几何朗兰兹对偶的一个解析版本（analytic version），其部分动因来自若干数学 \cite{L} 与物理 \cite{T} 的进展。这里考虑的不再是范畴以及作用在范畴上的函子，而是一个量子态希尔伯特空间，以及作用在该空间上的自伴算子，例如量子 Hitchin 哈密顿量。非常粗略地说，几何朗兰兹对偶的通常表述涉及 $G$-希格斯丛（Higgs bundle）模空间 $\M_H(G,C)$ 上全纯函数代数的形变量子化；而理论的解析版本涉及同一模空间（此时视为实辛流形）的普通量子化。照例，量子化意味着某一类适当的光滑函数——一般而言既非全纯也非反全纯——变成希尔伯特空间上的算子。关于量子 Hitchin 哈密顿量谱的规范场论解释，先前工作还可参见 \cite{Teschner:2010je,Balasubramanian:2017gxc,Nekrasov:2010ka,Nekrasov:2009rc,Nekrasov:2011bc,Nekrasov:2014yra,Jeong:2018qpc,Bonelli:2011na}。

\begin{yizhu}[译者按：两种量子化]
原文反复对比两种量子化，这是理解全文的一条主线。对复辛流形 $(\M_H,\Omega_I)$ 作\emph{形变量子化}（deformation quantization），是把全纯函数代数按 泊松括号形变成一个非交换代数；其表示论对应通常版本的几何朗兰兹。本文则把 $\M_H$ 的某个\emph{实}辛结构（例如 $\omega_J$）当作真正的辛形式，作普通（几何）量子化，得到希尔伯特空间与作用其上的自伴算子。二者的对象与结论都不同，不能混为一谈。参见 \cite{GuW}；相关记号见第 \ref{cm} 节。
\end{yizhu}

本文的目标是把几何朗兰兹对偶的解析版本放进规范场论框架。第一步只是要理解，在该框架中应当如何研究 $\M_H$ 作为实辛流形的量子化，以区别于把它作为复辛流形时的形变量子化。这里的基本想法是：实辛流形 $M$ 的量子化问题，是 $M$ 的某个适当复化 $Y$（如果这样的 $Y$ 存在的话）的 $A$-模型的一部分 \cite{GuW}。我们已在别处更详细地探讨过这一构造，作为本文的背景 \cite{GW}。应用到几何朗兰兹时，我们希望量子化希格斯丛模空间 $\M_H$，把它看作带有某个实辛结构的实辛流形。$\M_H$ 的一个适当复化，就是取两份 $\M_H$ 并赋予相反的复结构所得的乘积。以此作为理解 $\M_H$ 量子化的起点，我们将说明：把先前用来理解较传统版本几何朗兰兹的同一批规范场论（gauge theory）要素以新的方式组装起来，就能理解解析版本的几何朗兰兹。

本文组织如下。第 \ref{basic} 节解释用 $\N=4$ 超杨–米尔斯理论表述解析版几何朗兰兹对偶的基本设置。第 \ref{hitchval} 节解释电–磁对偶对 Hitchin 哈密顿量联合本征值的预言；第 \ref{wtw} 节分析 Hecke 算符或 't Hooft 算符以及对偶的 Wilson 算符。第 \ref{wkb} 节证明 Hitchin 哈密顿量的联合本征函数满足一个量子形变的 WKB 条件。第 \ref{realb} 节讨论希格斯丛模空间实形式的量子化。从第 \ref{avatar} 节开始，我们利用超对称边界条件交汇处出现的手征代数（chiral algebra），构造一类有趣的波函数，其谱分解由 $S$-对偶边界条件的几何控制。这些波函数通常不是本征函数，但我们预期它们在几何朗兰兹纲领解析版本的语境中会起有益作用。第 \ref{sec:chiral} 节通过把 Hecke 算符与 Kac--Moody 代数的谱流（spectral flow）自同构联系起来研究它们。Hecke 算符与手征代数谱流自同构之间的这一关系在此分析中很重要，先前已从不同但相关的角度被注意到 \cite{Teschner:2010je}。第 \ref{sec:symp} 节详细探讨一个具有显著性质的例子，它导致一个把 $SL_2$ 的三份 Hitchin 哈密顿量与 Hecke 算符交织起来的波函数。

若干进一步的问题在附录中处理。附录 \ref{coho} 分析沿典范余迷向膜（coisotropic brane）插入的局部算符的上同调。结果表明 Hitchin 的经典可积系统可以被量子化。附录 \ref{bmodel} 说明，电–磁对偶连同 $\M_H$ 量子化中希尔伯特空间内积的正定性，蕴含参数化全纯与反全纯 oper 的簇 $L_\op$ 与 $L_\bop$ 的交集必须是孤立且横截的（如 \cite{EFK} 所猜想，并在某些情形下已证明）。\zy{oper 是 Beilinson 与 Drinfeld 为几何朗兰兹引入的几何对象：粗略地说，它是带平坦联络的 $G$-丛，并带一个满足横截性条件的 Borel 约化。固定 $K_C^{1/2}$（等价地，固定 $C$ 上的旋量结构（spin structure））后，$\SL(2,\C)$ 的 oper 等价于 $C$ 上的射影结构（projective structure）；不同的 $K_C^{1/2}$ 选择给出不同的 oper 簇。参见 \cite{BD}。} 我们进一步指出，电–磁对偶蕴含 Hitchin 哈密顿量联合本征函数的一个自然归一化，并论证归一化波函数的希尔伯特空间范数大致由相应 oper 丛的挠率（torsion）给出。附录 \ref{somex} 探讨线算符所满足的一些微分方程的例子。附录 \ref{localmodel} 构造一个局部模型，用来分析 Hitchin 哈密顿量本征函数沿“不太稳定”丛（not very stable bundles）除子的奇异行为。

全文约定 $C$ 是亏格 $g>1$ 的黎曼曲面。所有讨论都可自然推广到带抛物结构的丛，但为简单起见我们省略这一推广（它已在 \cite{EFK,EFK2,EFK3} 中发展）。有足够多抛物点时，亏格 $g=0,1$ 也有相当类似的理论。没有抛物结构（或对 $g=0$ 抛物点太少）的 $g=0,1$ 情形需要不同处理，因为任何低能描述都需要规范场，而不只是 $\sigma$-模型场。
